\documentclass[aspectratio=169]{beamer}
\usepackage[T1]{fontenc}
\usepackage[default]{sourcesanspro}
\usepackage{newtxsf}
\usepackage{hyperref}
\usepackage{latexsym,amsmath,xcolor,multicol,booktabs,calligra}
\usepackage{graphicx,pstricks,listings,stackengine}
\graphicspath{{./figures/}}
\usepackage{lipsum}
\usepackage{wrapfig}
\usepackage{url}
\usepackage{subcaption}
\usepackage{listings}
\usepackage{minted}
\setminted{fontsize=\footnotesize}
\usepackage{fontawesome5}

\usepackage{etoolbox}

\author{Yan Lin}
\title{AI Systems \& Infrastructure}
\subtitle{Module B.3: Building AI-enabled Containers}
\institute{%
    \href{https://www.yanlincs.com}{www.yanlincs.com}
}
\date{October 26th, 2026}
\usepackage{AAU}
\AtBeginSubsection{}
\setbeamertemplate{navigation symbols}{}
\ifdefined\hidenotes\else
    \setbeameroption{show notes on second screen=right}
    \setbeamerfont{note page}{size=\large}
    \hypersetup{pdfkeywords={SP-Right}}
\fi

\def\cmd#1{\texttt{\color{red}\footnotesize $\backslash$#1}}
\def\env#1{\texttt{\color{blue}\footnotesize #1}}
\definecolor{aauyellow}{RGB}{167, 131, 55}
\definecolor{aaublue}{RGB}{32, 26, 78}
\definecolor{aaured}{RGB}{152, 43, 28}
\definecolor{textgrey}{RGB}{60, 60, 60}

\setbeamercolor{normal text}{fg=textgrey}
\usebeamercolor[fg]{normal text}

\def\emphasis#1{{\color{aaured}\textbf{#1}}}
\def\live#1{{\color{aaublue}\textbf{Live demo:}} #1}
\def\hr{\noindent\rule{\textwidth}{0.4pt}\\}
\renewcommand{\thefootnote}{}


\lstset{%
    basicstyle=\ttfamily\small,
    keywordstyle=\bfseries\color{deepblue},
    emphstyle=\ttfamily\color{deepred},
    stringstyle=\color{deepgreen},
    numbers=left,
    numberstyle=\small\color{halfgray},
    rulesepcolor=\color{red!20!green!20!blue!20},
    frame=shadowbox,
}

\begin{document}

\begin{frame}[plain,noframenumbering]
    \titlepage
    \vspace*{-0.6cm}
    \begin{figure}[htpb]
        \begin{center}
            \includegraphics[keepaspectratio, scale=0.5]{aau-logo-left-uk}
        \end{center}
    \end{figure}
\note{
}
\end{frame}


\begin{frame}
\tableofcontents[sectionstyle=show,
subsectionstyle=show/shaded/hide,
subsubsectionstyle=show/shaded/hide]
\note{
}
\end{frame}


\section{Introduction}

\begin{frame}{Recap of Module B.2: Container Fundamentals}{Container standard}
  Container as an industry standard for streamlining the process of deploying our AI systems onto any* computer.

  % TODO: Embed figures/module-b2-oci-standard.pdf

\note{
}
\end{frame}


\begin{frame}{Recap of Module B.2: Container Fundamentals}{Why containers deliver the premise}
  \begin{block}{Software deployment without containers}
    \begin{itemize}
      \item Setting up everything right on each computer the software is deployed to is tedious and prone to error
      \item Multiple pieces of software running on the same computer might interfere with each other
    \end{itemize}
  \end{block}

  \begin{block}{Software deployment with containers}
    \begin{itemize}
      \item A container \textbf{packs everything a piece of software needs in one self-contained package}
      \item A container provides an isolated environment for software to run without interference
    \end{itemize}
  \end{block}

\note{
}
\end{frame}


\begin{frame}{Recap of Module B.2: Container Fundamentals}{Use containers in practice}
  Docker ecosystem that includes all tools we need to work with containers.

  % TODO: Embed figures/module-b2-docker-ecosystem.pdf

\note{
}
\end{frame}


\begin{frame}{How to build a container image?}{The premise}
  Deploying our AI systems as containers will be easier and more robust than doing it manually on each computer.

  % TODO: Embed figures/module-b2-container-ship.pdf
\note{
}
\end{frame}


\begin{frame}{How to build a container image?}{The missing pieces}
  The prerequisite of deploying our AI systems as containers: having a container image packed with everything our AI system needs to run.

  \begin{itemize}
    \item Source code of our API server
    \item Source code and parameters of our AI model
    \item Configurations like API keys and rate-limiting thresholds
    \item Command to start the API server upon initialization of the container
  \end{itemize}

  All customized pieces that no existing container image on Docker Hub will have.
  
\note{
}
\end{frame}


\begin{frame}{Overview of this module}
  \begin{itemize}
    \item How to build a custom container image using a Dockerfile
    \item How to build and distribute our AI system as containers
  \end{itemize}

  % TODO: Embed figures/module-b2-docker-ecosystem-dockerfile.pdf
  
\note{
}
\end{frame}



\section{Building Images}

\begin{frame}[fragile]{The interactive approach, and why not}
    Drop into a container, install things, snapshot it.
    \begin{minted}{bash}
docker run -it --name builder python:3.13 bash
pip install fastapi uvicorn
exit
docker commit builder my-fastapi-image:dev
    \end{minted}

    Works, but rarely a good idea.
    \begin{itemize}
        \item No record of what we actually installed
        \item Nothing for teammates or our future selves to read
        \item Every change requires going through the same manual steps
        \item Throws away the layered structure that makes images efficient in the first place
    \end{itemize}
\note{
}
\end{frame}


\begin{frame}[fragile]{Dockerfile, the recommended way}
    Describe the image as code instead of building it interactively.
    A Dockerfile is a plain-text file that lists, step by step, what should go into the image.

    \begin{minted}{dockerfile}
FROM python:3.13-slim
WORKDIR /app
COPY hello.py .
CMD ["python", "hello.py"]
    \end{minted}

    Build with:
    \begin{minted}{bash}
docker build -t my-hello:1.0 .
    \end{minted}

    The \texttt{.} is the build context: the set of files Docker has access to during the build.
\note{
}
\end{frame}


\begin{frame}{Common Dockerfile instructions}
    \begin{itemize}
        \item \texttt{FROM <image>}: pick the base image, always first
        \item \texttt{WORKDIR <path>}: set the working directory inside the image
        \item \texttt{COPY <src> <dest>}: bring files in from the build context
        \item \texttt{RUN <cmd>}: execute a shell command during the build
        \item \texttt{ENV KEY=value}: set an environment variable visible at runtime
        \item \texttt{EXPOSE <port>}: document the port, does not actually publish it
        \item \texttt{CMD ["program", "arg"]}: default command when a container starts
    \end{itemize}
\note{
}
\end{frame}


\begin{frame}[fragile]{Layer order matters}
    A Dockerfile maps closely to the layers of the resulting image, roughly one layer per instruction.
    Docker caches layers between builds.
    As soon as one layer's input changes, every layer on top of it has to be rebuilt as well.

    \begin{block}{Worse order}
        \begin{minted}{dockerfile}
COPY . .
RUN pip install -r requirements.txt
        \end{minted}
        Every edit to \texttt{main.py} invalidates the \texttt{COPY}, which forces \texttt{pip install} to rerun.
    \end{block}

    \begin{block}{Better order}
        \begin{minted}{dockerfile}
COPY requirements.txt .
RUN pip install --no-cache-dir -r requirements.txt
COPY . .
        \end{minted}
        \texttt{pip install} only reruns when \texttt{requirements.txt} changes.
    \end{block}
\note{
}
\end{frame}



\section{Containerizing the AI Server}

\begin{frame}{Project layout}
    The end goal is a single image that bundles our code, its Python dependencies, and the model weights, so that anyone with our image can bring up the server with one \texttt{docker run}.

    Files we need to containerize the AI API server from Module A.4:
    \begin{itemize}
        \item \texttt{main.py} and other Python files of the server
        \item \texttt{requirements.txt} listing the Python packages we need
        \item the \texttt{Dockerfile} we are about to write
    \end{itemize}
\note{
}
\end{frame}


\begin{frame}[fragile]{A first Dockerfile}
    \begin{minted}{dockerfile}
FROM python:3.13-slim

WORKDIR /app

COPY requirements.txt .
RUN pip install --no-cache-dir -r requirements.txt

COPY main.py ./

EXPOSE 8000

CMD ["uvicorn", "main:app", "--host", "0.0.0.0", "--port", "8000"]
    \end{minted}

    Note the \texttt{--host 0.0.0.0} argument.
    By default Uvicorn binds to \texttt{127.0.0.1}, which inside a container only accepts traffic from that container.
    Binding to \texttt{0.0.0.0} makes it listen on all interfaces, so port mapping can actually reach it.
\note{
}
\end{frame}


\begin{frame}[fragile]{Build and run}
Build the container image.
    \begin{minted}{bash}
docker build -t my-ai-server:1.0 .
    \end{minted}
    First build is slow because PyTorch is large.
    Subsequent builds are fast as long as \texttt{requirements.txt} does not change.

Then run it with port mapping.
    \begin{minted}{bash}
docker run -d --name ai-server -p 8000:8000 my-ai-server:1.0
    \end{minted}

    Now \texttt{http://127.0.0.1:8000/v1/classify} should be reachable from the A.2 client, just like in A.4.
    The difference: everything the server needs lives inside the container and is ready to be distributed at scale.
\note{
}
\end{frame}


\begin{frame}[fragile]{Push to a registry}
    A built image lives only on the machine where it was built until we push it to a registry.
    Like Git repositories, container images are pushed and pulled to and from registries.

    \begin{minted}{bash}
# 1. Log in (Docker Hub example)
docker login

# 2. Tag the local image with the registry path
docker tag my-ai-server:1.0 yourusername/my-ai-server:1.0

# 3. Push
docker push yourusername/my-ai-server:1.0
    \end{minted}

    Tag format: \texttt{[<registry>/]<namespace>/<repository>:<tag>}.
    For Docker Hub, the registry part is omitted.
    Anyone can fetch your image using \texttt{docker pull yourusername/my-ai-server:1.0}.
\note{
}
\end{frame}



\section{Layers vs Volumes}

\begin{frame}{The recurring design choice}
    For each file, we have to decide: inside the image as a layer, or mounted in from the outside as a volume?

    \begin{block}{General principle}
\begin{itemize}
\item \textbf{Inside:} things that should be the same across every running instance
\item \textbf{Volume:} things that should differ between instances and/or persist beyond one instance
\end{itemize}
    \end{block}

    \begin{block}{Examples}
    \begin{itemize}
        \item Application code and Python packages: inside, since every instance runs the same
        \item User database: volume, differ between instances and should survive container restart
        \item Model weights: can be inside or volume or even non-persistent, each with caveats
    \end{itemize}
    \end{block}
\note{
}
\end{frame}


\begin{frame}[fragile]{Baking model weights into the image}
    For small models, download at build time is fine:
    \begin{minted}{dockerfile}
COPY requirements.txt .
RUN pip install --no-cache-dir -r requirements.txt

# Pre-download ResNet-18 weights into a layer
RUN python -c "from torchvision.models import \
                   resnet18, ResNet18_Weights; \
               resnet18(weights=ResNet18_Weights.DEFAULT)"

COPY main.py db.py ./
    \end{minted}

    A fresh container now starts immediately, with no runtime download.
    For very large models, consider falling back to downloading at startup.
\note{
}
\end{frame}


\begin{frame}[fragile]{Persisting the database via a volume}
    The SQLite database file stores our users, API keys, and request logs, and it must survive container removal.
    If the database lives inside the writable layer, removing the container deletes all our data.

    \begin{minted}{bash}
mkdir -p data

docker run -d --name ai-server \
    -p 8000:8000 \
    -v $(pwd)/data:/app/data \
    -e DATABASE_URL=sqlite:////app/data/ai_api.db \
    my-ai-server:1.0
    \end{minted}

    \begin{itemize}
        \item Database lives on the host, survives container removal
        \item Server reads \texttt{DATABASE\_URL} from the environment
        \item We can swap the container for a newer image and keep the data
    \end{itemize}
\note{
}
\end{frame}



\section{GPU Passthrough}

\begin{frame}{Containers do not see the GPU by default}
    If we run our server on a machine with an NVIDIA GPU, the container needs explicit access to the GPU device.
    Containers are isolated from host hardware, and the GPU has to be passed through.

    \begin{itemize}
        \item \texttt{--device <host>} grants access to one specific device file
        \item Same pattern works for webcams, serial ports, USB, sound cards
        \item \texttt{--privileged} drops all device restrictions at once, but strips away most of the isolation containers exist to provide
    \end{itemize}

    GPUs need more than raw device files. They also need userspace driver libraries that match the kernel module on the host.
\note{
}
\end{frame}


\begin{frame}[fragile]{NVIDIA Container Toolkit}
    The standard tool for NVIDIA GPUs.
    Handles the userspace libraries and bridges the host driver to the container.

    \begin{minted}{bash}
docker run -d --name ai-server \
    --gpus all \
    -p 8000:8000 \
    my-ai-server:1.0
    \end{minted}

    \begin{itemize}
        \item \texttt{--gpus all}: every GPU on the host
        \item \texttt{--gpus "device=0"}: only one specific GPU
        \item \texttt{--gpus 2}: limit to two GPUs
    \end{itemize}
\note{
}
\end{frame}


\begin{frame}{Other AI hardware}
    The \texttt{--gpus} flag only covers NVIDIA.

    \begin{itemize}
        \item AMD GPUs use ROCm. Access via \texttt{--device /dev/kfd /dev/dri}, or AMD's container toolkit
        \item Google TPUs require some more complicated setup
        \item Consumer NPUs (Apple Neural Engine, Qualcomm, AMD Ryzen AI) mostly are not exposed to containers
    \end{itemize}
\note{
}
\end{frame}



\section{Exercise}

\begin{frame}{Containerize the A.4 server}
    Turn the API server from the Module A.4 exercise into a container image.
    Point the A.2 client at it and confirm end-to-end behavior is unchanged.

    \begin{block}{Outline}
        \begin{enumerate}
            \item Write a \texttt{Dockerfile} to containerize the API server
            \item Run \texttt{docker build} with a tag you can remember. Watch the layers in the build log
            \item Run the container and verify the A.2 client can reach it
        \end{enumerate}
    \end{block}
\note{
}
\end{frame}


\begin{frame}{Extensions}
    \begin{enumerate}
        \item Pre-download the model weights at build time. Confirm with \texttt{docker logs} that no download happens at runtime
        \item Mount a host directory for the SQLite database. Register a user, stop and remove the container, start a new one with the same volume, verify data survived
        \item If you have an NVIDIA GPU, install the toolkit, run with \texttt{--gpus all}, move the model to \texttt{cuda}
        \item Push the image to Docker Hub. Pull on a different machine and confirm one \texttt{docker run} brings the server up
    \end{enumerate}
\note{
}
\end{frame}


\begin{frame}{What you should walk away with}
    By the end of the exercise, you should be able to:
    \begin{enumerate}
        \item Build a container image with a Dockerfile
        \item Containerize an AI API server, including its model and GPU access
        \item Decide what belongs in the image versus a mounted volume
    \end{enumerate}

    \vspace{0.5em}
    Our AI server now ships as a single image that runs the same anywhere.
    In the next module, we rent a computer in the cloud and run it there.
\note{
}
\end{frame}



\appendix

\begin{frame}[plain,noframenumbering]
\begin{center}
{\large Thank you!}
\vspace{1cm}
\\[0.5em]
{\small \faEnvelope\ \href{mailto:lyan@cs.aau.dk}{lyan@cs.aau.dk} \\
\faGlobe\ \href{https://www.yanlincs.com}{www.yanlincs.com}}
\end{center}
\note{
}
\end{frame}

\end{document}
